\documentclass[11pt]{article}
\usepackage[margin=1in]{geometry}
\usepackage{amsmath,amssymb}
\usepackage{tikz}
\usepackage{float}
\usepackage{graphicx}
\usepackage[hidelinks]{hyperref}

\setlength{\parindent}{0pt}
\setlength{\parskip}{0.7\baselineskip}

\begin{document}

\section{Local and Heuristic-Non-Optimal Searches}

So far we've talked about searches that are guaranteed to find an optimal answer (or the correct answer if decision version).

This week we're talking about other kinds of search which in general are fast but not guaranteed to find best-possible answer.

\subsection{Local search}

The `local' in local search refers to a local exploration of a search space.

Think about all the possible solutions to one of our problems: e.g.\ satisfiability problems. Here a possible solution is a set of truth assignments to variables: then a local search might look around `similar' assignments and gradually try to move toward an overall satisfying assignment. We need some notion of a `neighbour' state, and some indication of what is a good or a bad state to help us direct the search.

If we're talking about truth assignments:
\begin{itemize}
    \item a neighbour truth assignment might be one with a small number of different assignments, and
    \item a measure of goodness might be how many clauses are satisfied.
\end{itemize}

We'll talk about:
\begin{itemize}
    \item hill climbing
    \begin{itemize}
        \item we'll mention exhaustive local search at the end if we have time
    \end{itemize}
    \item simulated annealing
    \item ant colony algorithms
    \item Tabu search
    \item genetic algorithms
\end{itemize}

\subsection{Hill climbing}

Hill climbing is based on the idea that the solution space is like a landscape with better solutions higher up. Hill climbing is like trying to ascend a hill by always stepping only upwards.

Example: live example with vertex cover.

Let's look at some pseudocode (courtesy of Wikipedia):

\begin{figure}[H]
    \centering
    \includegraphics[width=0.7\textwidth]{image.png}
    \caption{Discrete-space hill climbing (from Wikipedia).}
\end{figure}

An important observation: hill climbing will eventually give an optimal answer for a \textbf{convex problem} -- this is a problem without local maxima. (I'll draw an example squiggly line.)

Hill climbing is good at this kind of thing:

\begin{figure}[H]
    \centering
    \includegraphics[width=0.6\textwidth]{image-1.png}
    \caption{A relatively smooth hill with one peak.}
\end{figure}

But not great at these kinds of things:

\begin{figure}[H]
    \centering
    \includegraphics[width=0.6\textwidth]{image-2.png}
    \caption{A jagged range of mountains with several peaks.}
\end{figure}

Do we think that satisfiability is a convex problem? We'll try to concoct an example. Dependent on the `neighbour' relation and the `goodness' measure.

Hill climbing variants (we'll look at little 1D cartoon examples of these):
\begin{itemize}
    \item steepest-ascent
    \begin{itemize}
        \item instead of going to the first `higher' neighbour, find the most-better neighbour and go there.
    \end{itemize}
    \item stochastic
    \begin{itemize}
        \item move to a random neighbour with probability related to how much better than neighbour is
    \end{itemize}
    \item shotgun (random-restart)
    \begin{itemize}
        \item do a bunch of hill climbing, restarting occasionally. Keep the best result.
    \end{itemize}
\end{itemize}

\subsection{Simulated annealing}

The name comes from metallurgy: heating and cooling a metal to change its characteristics. It's not a lot like it really, but the `cooling' is the analogy.

Intuition:
\begin{itemize}
    \item every state has some `energy' (and here we look for low-energy states)
    \item the system has a `temperature' that determines how easy it is to jump around the system
    \begin{itemize}
        \item cooler means less jumping
    \end{itemize}
    \item we gradually `cool' the system as we continue our search
    \item we quit when an answer is good enough, or we're not moving, or we run out of computation
\end{itemize}

Let's look at an intuition-giving animation:

{\small\url{https://en.wikipedia.org/wiki/File:Hill_Climbing_with_Simulated_Annealing.gif}}

Let's look at some pseudocode:

\begin{verbatim}
P(energy_1, energy_2, temperature) is the probability of acceptance function
E(s) is the energy of state s

s <-- s_0
for k in [0 .. k_{max}]:
    T <-- temperature schedule at time k
    s_{new} <-- a random neighbour of s
    if P(E(s), E(s_{new}, T)) >= random number in (0, 1):
        s <-- s_{new}
return s
\end{verbatim}

Setting the $P$ and $E$ functions, temperature schedule, as well as the parameters $k_{max}$ is problem-specific. Sometimes the neighbour generation function includes longer jumps as well (helps with escaping local maxima/minima).

Simulated annealing is popular for travelling salesperson (TSP).

A nice set of examples and visualisations here:

{\small\url{https://www.fourmilab.ch/documents/travelling/anneal/}}

\subsection{Tabu search}

Called `tabu' because we build up sets of things we `don't touch'. Really it's a family of approaches.

Intuition:
\begin{itemize}
    \item as with other local searches we generate neighbour states
    \item build up `banned' lists of moves that we do not want to consider
    \begin{itemize}
        \item prevents cycling
    \end{itemize}
    \item very related to simulated annealing
\end{itemize}

Let's look at pseudocode (courtesy of Wikipedia):

\begin{figure}[H]
    \centering
    \includegraphics[width=0.75\textwidth]{image-3.png}
    \caption{Pseudocode for tabu search (from Wikipedia).}
\end{figure}

We can get fancier with other kinds of `memory' lists -- but that's a bit beyond what we'll talk about.

Also suitable for TSP.

\subsection{Ant colony algorithms}

Based on the behaviour of ants when searching for food -- specifically for problems that can be phrased as finding a path through a graph. Simulated ants wander around a parameter/solution space looking for solutions. The `ants' lay down little markers indicating where they've found good solutions. The markers `evaporate' over time, so shortest paths have more pheromone and attract more ants -- feedback cycle ensues.

Some (very) pseudocode:

\begin{verbatim}
while we're not done:
    put each ant on a starting node
    while an ant hasn't found a solution:
       for each ant:
         move ant to next node using transition rule
         apply the pheromone update
    update the best solution
    do the secondary pheromone update
report the best solution
\end{verbatim}

Transition rules and pheromone updates are problem-specific.

\subsection{Genetic algorithms}

Supposedly inspired by evolution in the natural world.

Intuition: given a solution, use genetic operations to generate better solutions.

We need:
\begin{itemize}
    \item some genetic-style representation
    \item a fitness function
    \item definitions of how we do our mutation, crossover, selection
\end{itemize}

We maintain a `population' of solutions. Then we select for the best by some measure. Then we mutate or combine them, then we repeat.

\end{document}
